\section{The arithmetic positivity target}\label{sec:arithmetic}

This section fixes the target independently of any field-theory model.
The formulas also explain why positivity of a single covariance kernel
would be insufficient.

\subsection{The compactly supported Weil form}

Let $I_L=(-L/2,L/2)$, with $L>0$, and let
$f\in C_c^\infty(I_L;\C)$. Write $F$ for its zero extension to $\R$ and use
\begin{equation}
 \widehat F(\tau)=\int_\R e^{-\ii\tau x}F(x)\dd x,
 \qquad \ip f g=\int_{I_L}\overline{f(x)}g(x)\dd x.
\end{equation}
For $a>0$, define the compressed shift $T_af(x)=F(x-a)$ on $I_L$.
Let $\psi=\Gamma'/\Gamma$ be the digamma function, and let
$\Lambda(n)=\log p$ when $n=p^k$, $p$ prime and $k\geq1$, and zero otherwise.
Finally put
\begin{equation}
 C(f)=\int_{I_L}f(x)\cosh(x/2)\dd x,
 \qquad S(f)=\int_{I_L}f(x)\sinh(x/2)\dd x.
\end{equation}
The normalization of the program is
\begin{align}\label{eq:weil}
 Q_{0,L}[f]={}&\frac1{2\pi}\int_\R
 \bigl[\Ree\psi(\tfrac14+\tfrac{\ii\tau}2)-\log\pi\bigr]
 |\widehat F(\tau)|^2\dd\tau
 +2|C(f)|^2-2|S(f)|^2\notag\\
 &-\sum_{\substack{n\geq2\\\log n<L}}
 \frac{\Lambda(n)}{\sqrt n}\ip f{(T_{\log n}+T_{\log n}^*)f}.
\end{align}
The prime-power sum is finite on a fixed interval. The pole term is
indefinite, and the archimedean multiplier includes its local constant.
Neither may be dropped in a positivity argument.

For comparison with the usual explicit formula, set
$h=F*\widetilde F$, where $\widetilde F(x)=\overline{F(-x)}$.
The two prime sums in the Weil functional give the last line of
\eqref{eq:weil}; its $e^{x/2}+e^{-x/2}$ term gives
$2|C|^2-2|S|^2$. Its archimedean distribution has Fourier multiplier
$\Ree\psi(1/4+\ii\tau/2)-\log\pi$. Thus
$Q_{0,L}[f]=W(F*\widetilde F)$ in the convention of
\cite[p.~1]{Suzuki}. In particular the weight is $\Lambda(n)/\sqrt n$.

\begin{theorem}[Weil positivity criterion, used as background]\label{thm:weil}
RH holds if and only if $Q_{0,L}[f]\geq0$ for every $L>0$ and every
$f\in C_c^\infty(I_L;\C)$.
\end{theorem}
The equivalence is the classical criterion, not a result of the present
construction. Every compact support lies in one such interval, so an
unbounded sequence of interval lengths suffices. A finite collection of
lengths or numerical matrices does not.

\subsection{The positive kinetic term and the indispensable local term}

Put $a_n=2n+1/2$, for $n\geq0$, and
\begin{equation}\label{eq:b-symbol}
 b(\tau^2)=\Ree\psi(\tfrac14+\tfrac{\ii\tau}2)-\psi(\tfrac14)
 =2\sum_{n=0}^\infty\frac{\tau^2}{a_n(a_n^2+\tau^2)}\geq0,
 \qquad w_0=\psi(\tfrac14)-\log\pi<0.
\end{equation}
The constant $\psi(1/4)$ is not equal to $\log\pi$. Adding and
subtracting $\psi(1/4)$ in the original multiplier gives the exact identity
\begin{align}\label{eq:archimedean-split}
 \Ree\psi(\tfrac14+\tfrac{\ii\tau}2)-\log\pi
 &=[\Ree\psi(\tfrac14+\tfrac{\ii\tau}2)-\psi(\tfrac14)]
   +[\psi(\tfrac14)-\log\pi]\notag\\
 &=b(\tau^2)+w_0.
\end{align}
Thus $w_0$ is the remainder after separating the nonnegative multiplier
$b(\tau^2)$; it is already present in \eqref{eq:weil}, not an additional
term introduced into the form.

The series in \eqref{eq:b-symbol} follows by subtracting the convergent digamma partial
fractions. Each exponential in \eqref{eq:intro-tower} satisfies
$\int_0^\infty e^{-a u}(1-\cos\tau u)\dd u
=\tau^2/[a(a^2+\tau^2)]$. Plancherel therefore gives
\begin{equation}\label{eq:kinetic}
 \mathcal E_\Gamma[F]
 :=\int_0^\infty\nGamma(u)\norm{F(\,\cdot+u)-F}_2^2\dd u
 =\frac1{2\pi}\int_\R b(\tau^2)|\widehat F(\tau)|^2\dd\tau.
\end{equation}
Using \eqref{eq:archimedean-split} and
$\frac1{2\pi}\int_\R|\widehat F(\tau)|^2\dd\tau=\norm f_2^2$,
the archimedean contribution becomes
\begin{equation}\label{eq:archimedean-energy}
 \frac1{2\pi}\int_\R
 [\Ree\psi(\tfrac14+\tfrac{\ii\tau}2)-\log\pi]
 |\widehat F(\tau)|^2\dd\tau
 =\mathcal E_\Gamma[F]+w_0\norm f_2^2.
\end{equation}
Substituting this identity into \eqref{eq:weil} rewrites the full target as
\begin{align}\label{eq:weil-energy}
 Q_{0,L}[f]={}&\mathcal E_\Gamma[F]+w_0\norm f_2^2
 +2|C(f)|^2-2|S(f)|^2\notag\\
 &-\sum_{\substack{n\geq2\\\log n<L}}
 \frac{\Lambda(n)}{\sqrt n}\ip f{(T_{\log n}+T_{\log n}^*)f}.
\end{align}
Here $w_0\norm f_2^2=w_0\int_{I_L}|f(x)|^2\dd x$ is the fixed
negative local term. The positive difference energy is only one part of
the desired form: its local term, pole contribution and prime-power
translations must all be retained.

The natural form domain on a fixed interval is given by zero extensions
with $\int\log(2+|\tau|)|\widehat F(\tau)|^2\dd\tau<\infty$.
Indeed $b(\tau^2)$ grows logarithmically at infinity, whereas the pole
and finite translation terms are bounded forms on $L^2(I_L)$. This
explains the logarithmic form domain and semiboundedness of the
localized form. For the logical proof target in Theorem~\ref{thm:weil},
the smooth compactly supported class is enough.

\subsection{Shifted transfers and a sufficient contraction route}

Define the entire completed zeta function and its centered version by
\begin{equation}
 \xi(s)=\tfrac12s(s-1)\pi^{-s/2}\Gamma(s/2)\zeta(s),
 \qquad \Xi(p)=\xi(\tfrac12+p).
\end{equation}
For $0<\omega\leq1/2$, the shifted transfer is
\begin{equation}\label{eq:transfer}
 K_\omega(p)=\frac{\Xi(p-\omega)}{\Xi(p+\omega)}.
\end{equation}
The functional equation and reality imply $|K_\omega(\ii\tau)|=1$
where the ratio is defined. This boundary identity is unconditional;
it is not a proof that an unweighted convolution operator is a
contraction. In particular poles away from that boundary cannot be
ignored.

The causal interpretation uses inverse Laplace transform on a line
$\Ree p=\sigma>1$. Let $k_\omega$ be that causal distribution and
$V_{\omega,L}=P_L(k_\omega*)P_L$, where $P_L$ denotes restriction to
$I_L$ together with zero extension when used as an input. The explicit
kernel below is locally integrable for $\omega>0$, so these finite-window
operators are bounded. At $\omega=0$, $k_0=\delta_0$ and $V_{0,L}=I$.
Logarithmic differentiation gives
\begin{equation}
 \partial_\omega K_\omega(p)=-a_\omega(p)K_\omega(p),\qquad
 a_\omega(p)=\frac{\Xi'}{\Xi}(p-\omega)
             +\frac{\Xi'}{\Xi}(p+\omega).
\end{equation}
At zero shift the causal generator has Laplace symbol
\begin{align}\label{eq:generator}
 a_0(p)={}&\frac2{p+1/2}+\frac2{p-1/2}-\log\pi
 +\psi(\tfrac14+\tfrac p2)
 -2\sum_{n\geq2}\frac{\Lambda(n)}{\sqrt n}e^{-p\log n}.
\end{align}
The identity
\begin{equation}
 \psi(\tfrac14+\tfrac p2)-\psi(\tfrac14)
 =2\int_0^\infty\nGamma(u)(1-e^{-pu})\dd u
\end{equation}
shows that the real quadratic form of this causal generator is exactly
\eqref{eq:weil}. The two rational terms give the symmetric kernel
$2\cosh((x-y)/2)$; its form is $2|C|^2-2|S|^2$.
On smooth test functions the first variation is therefore
\begin{equation}\label{eq:first-variation-arithmetic}
 \ip f{(I-V_{\omega,L}^*V_{\omega,L})f}
 =2\omega Q_{0,L}[f]+o(\omega),\qquad\omega\downarrow0.
\end{equation}
One obtains this expansion by differentiating the causal kernel at
$\omega=0$ using \eqref{eq:generator}; smooth zero extensions make the
subtracted difference integral convergent in $L^2$ on a finite window.
Thus an independent proof of $\norm{V_{\omega,L}}\leq1$ for every $L$
and shifts tending to zero would imply RH. We need only this sufficient
implication here; no global unweighted realization of $K_\omega$ is
assumed.

\subsection{The exact positive part, and what it leaves unresolved}

Write $a=1/2-\omega$, $b=1/2+\omega$, $s=1/2+p$, and
$\mathcal L(s)=\pi^{-s/2}\Gamma(s/2)\zeta(s)$, distinguished from the
von Mangoldt function. Direct factorization yields
\begin{equation}\label{eq:factor}
 K_\omega(p)=\frac{(p+a)(p-b)}{(p+b)(p-a)}\,
 \widetilde K_\omega(p),\qquad
 \widetilde K_\omega(p)=\frac{\mathcal L(s-\omega)}{\mathcal L(s+\omega)}.
\end{equation}
The beta integral and Euler product, for $\Ree p>b$, give
\begin{align}
 g_\omega(t)&=\frac{2\pi^\omega}{\Gamma(\omega)}
 e^{-a t}(1-e^{-2t})^{\omega-1}\mathbf1_{t>0},\label{eq:beta}\\
 c_n(\omega)&=n^{\omega-1/2}\prod_{\ell\mid n,\ \ell\text{ prime}}
 (1-\ell^{-2\omega}),\qquad c_1=1,\label{eq:comb}\\
 \widetilde k_\omega&=g_\omega*\sum_{n\geq1}c_n(\omega)\delta_{\log n}\geq0.
 \label{eq:positive-part}
\end{align}
For example, the substitution $r=e^{-2t}$ in the Laplace transform of
\eqref{eq:beta} gives
$\pi^\omega\Gamma((p+a)/2)/\Gamma((p+b)/2)$.
For the zeta quotient the coefficient at a prime power $\ell^k$ is
$\ell^{k\omega}(1-\ell^{-2\omega})$ before multiplication by
$n^{-1/2}$, proving \eqref{eq:comb} by multiplicativity.
Only finitely many shifts contribute on a fixed window. These formulas
define positive locally finite measures, not normalized probability
laws; the endpoint $a=0$ in particular precludes an unqualified
Beta$(a/2,\omega)$ probability interpretation.

Absorb one rational factor into the positive part:
\begin{align}
 \widehat k_\omega
 &=\widetilde k_\omega+2a(e^{a\,\cdot}\mathbf1_{>0})*
 \widetilde k_\omega\geq0,\notag\\
 k_\omega&=\widehat k_\omega
 -2b(e^{-b\,\cdot}\mathbf1_{>0})*\widehat k_\omega.
 \label{eq:rational-subtraction}
\end{align}
Thus the full kernel is a positive part minus an exponential average.
This exact reduction was central to the parent Loewner
investigation~\cite{ParentLoewner}. It neither makes $k_\omega$ positive
nor bounds the norm of its finite-window convolution. The required
cancellation is still the arithmetic problem.

The fractional-dimension interpretation follows from the radial beta
integral
\begin{equation}\label{eq:fractional}
 \int_{\R^d}(1+|u|^2)^{-(p+b)/2}\dd^d u
 =\pi^{d/2}\frac{\Gamma((p+b-d)/2)}{\Gamma((p+b)/2)}.
\end{equation}
For integer $d$ this is an ordinary integral in its convergence region;
its radial continuation at $d=2\omega$ reproduces the gamma factor.
It is not a continuously variable lattice dimension. To see a concrete
obstruction to one natural lattice continuation, put
$\vartheta(q)=\sum_{k\in\mathbb Z}q^{k^2}$ and $m=1+d$.
The coefficient of $q^3$ in $\vartheta(q)^m$ is
\begin{equation}
 8\binom m3=\tfrac43m(m-1)(m-2)=\tfrac43(d^3-d)<0
 \quad(0<d<1).
\end{equation}
It cannot be a positive point count. This excludes that interpolation,
not every positive kernel or analytic continuation. Neither
\eqref{eq:fractional} nor a bare Loewner map supplies the prime comb and
rational subtraction in \eqref{eq:positive-part}--\eqref{eq:rational-subtraction}.
